\subsection{Three exact local lifts of each normalized cubic state}

\paragraph{D23.H3 (canonical detector and its local lifts).}
Let $A,B,C\in\mathfrak S_{1344}$ be D23.P17's literal braid
generators, and let $V$ be the actual full-GT conjugator of
D23.P18. Suppose $\rho=0\pmod3$. On each of the four
orthogonal blocks numbered $7,8,10,11$, the restriction of
$V$ is one of exactly three explicitly recorded projective
linear maps, for either value $a_3\in\{0,1\}$.
All these maps have determinant one, quadratic multiplier
one, and spinor norm one. The candidates for the two values
of $a_3$ need not be the same maps.

The degree-$36$ quotient used to impose this constraint is
intrinsic to the finite braid action. In particular the
argument does not assume that deletion of a tuple coordinate
commutes with the full-GT action on tuples.

\begin{proof}
The group $G_3=\langle A,B\rangle$ has the following orbit
degrees, counted with multiplicity:
\[
                 4,8,8,12,80,108,160,160,240,240,324.
\]
The detector subset of D23.P17 is exactly the union of the
unique orbits of degrees $108$ and $324$. The Drinfeld action
preserves the three-strand subgroup and is an automorphism
there. Thus $V$ normalizes $G_3$ and preserves both of these
orbits separately.

On the degree-$108$ orbit the permutation image has order
$11664$ and exactly one minimal nontrivial block system.
It has $36$ blocks of size three and agrees with the
restriction of D23.P17's detector map. On the degree-$324$
orbit the image again has order $11664$. There are exactly
four minimal nontrivial block systems, each with $108$ blocks
of size three. The join of these four equivalence relations
has $36$ classes of size nine, and is again exactly the
detector map. A normalizer permutes the minimal block
systems, so it preserves their join. Both degree-$36$
quotients are therefore preserved intrinsically.

Their labelled braid generators are the same literal
degree-$36$ generators in D23.H1. The induced permutations
must satisfy that finding's simultaneous conjugacy equations.
Its uniqueness assertion identifies both induced permutations
with the same recorded $\theta_{0,a_3}$. Hence, writing $\psi$
for the detector map on its $432$-point domain, we have
\[
                      \psi(Vv)=\theta_{0,a_3}\psi(v).
\tag{D23.H3.1}
\]
This uses only a canonical finite quotient and the actual
braid covariance, not a topological forgetting assertion.

The twelve blocks are permuted by $V$, by D23.F30's socle
description. Alternatively the unique minimal block system
of the full transitive action gives this permutation directly.
The unique degree-$108$ orbit meets precisely blocks
$7,8,10,11$, in $27$ points apiece. Their detector images are
four disjoint sets of nine labels. Both
$\theta_{0,0}$ and $\theta_{0,1}$ preserve each of these four
sets. Equation (D23.H3.1) consequently fixes those four
blocks individually. The normalizer assertion D23.C24
places each local map in the projective similitude group of
the recorded quadratic form. The same conclusion also
follows from the cohomological construction D23.P19, which
already restricts the multiplier to one.

Use the group-transversal identifications of each block with
block zero. In each of these four blocks the full detector
domain has $108$ points; every one of the $36$ labels has
exactly three preimages. Choose the first six linearly
independent projective coordinate vectors in that domain.
They form a basis of the six-dimensional vector space.
By (D23.H3.1), each basis vector has three possible projective
images. Fixing the scalar of the first image leaves two
choices of scalar for each of the other five. This exhausts
all possible projective linear maps satisfying (D23.H3.1)
on the chosen basis: there are at most $3^6 2^5$ candidates.

For each multiplier $\mu\in\{1,2\}$, the certificate first
imposes the pairwise Gram equations on the basis images.
Their diagonal equations vanish because all six basis
vectors and all allowed images are singular. Exact rank
checks discard singular maps. Exactly six projective
similitudes remain at this stage in each of the eight cases
$(a_3,j)$, $a_3=0,1$, $j=7,8,10,11$. Testing (D23.H3.1)
on all $108$ points leaves exactly three in every case.
Every surviving map has $(\det,\mu)=(1,1)$.
The independent reflection certificate computes spinor
norm one for each. No possible projective similitude was
omitted: its six images and their relative scalars occur
in precisely this enumeration.

The complete permutation and matrix lists, together with
the canonical block-system checks, are in
\path{records/hurwitz_h1_orthogonal_lifts.npz} and the JSON
record of the same name. The executable certificate is
\path{code/hurwitz_h1_orthogonal_lifts.py}. Independent
checks of the fixed blocks and all spinor norms are in
\path{agents/characters/records/hurwitz_h3_independent_spinor.json}.
This establishes necessary local candidates, not the
existence of an actual GT element for any of them.
\end{proof}
